% Generated from locale/tr/content/counterfactuals/introduction/material-conditional.tex; do not hand-edit.


\olfileid[tr]{cnt}{int}{mat}

\olsection{Maddi Koşul}

Bir koşullu tümcenin İngilizcedeki en yalın biçimi ``Eğer \dots{} ise
\dots'' kalıbındadır; burada \dots{} ile gösterilen parçaların kendileri de
tümcedir. Örneğin: ``Eğer bunu uşak yaptıysa bahçıvan masumdur.'' Mantığa
giriş derslerinde koşullu tümceleri $\lif$ bağlacıyla simgeleştirmeyi
öğreniriz: \dots{} ile belirtilen parçaları örneğin $!A$ ve~$!B$
!!{formula}s{}iyle simgeler, bütün koşullu tümceyi de $!A\lif !B$ olarak
gösteririz.

$\lif$ bağlacı \emph{doğruluk fonksiyonludur}; yani $!A\lif !B$'nin
doğruluk değeri---$\True$ ya da $\False$---$!A$ ile~$!B$'nin doğruluk
değerlerince belirlenir: $!A\lif !B$ ancak ve ancak $!A$ yanlış ya da $!B$
doğruysa doğrudur; aksi hâlde yanlıştır. Bir doğruluk değeri ataması
$\pAssign v$'ye göre $\pSat{v}{!A\lif !B}$ ancak ve ancak
$\pSat/{v}{!A}$ ya da $\pSat{v}{!B}$ olacak biçimde tanımlanır. Bu
semantiğe sahip $\lif$ bağlacına \emph{maddi koşul} denir.

Bu tanımdan birtakım temel mantıksal olgular çıkar. İlk olarak, maddi koşul
için tümdengelim teoremi geçerlidir:
\begin{align}
  \text{$\Gamma,!A\Entails !B$ ise } \Gamma\Entails !A\lif !B
\end{align}
Doğruluk fonksiyonludur: $!A\lif !B$ ile $\lnot !A\lor !B$ denktir:
\begin{align}
  !A \lif !B & \Entails \lnot !A \lor !B\\
  \lnot !A \lor !B & \Entails !A \lif !B
  \intertext{Bir maddi koşul hem artbileşeninden hem önbileşeninin
    değillemesinden çıkar:}
  !B & \Entails !A \lif !B\\
  \lnot !A & \Entails !A \lif !B
  \intertext{Yanlış bir maddi koşul, önbileşeniyle artbileşeninin
    değillemesinin tümel evetlemesine denktir: $!A\lif !B$ yanlışsa
    $!A\land\lnot !B$ doğrudur ve bunun tersi de geçerlidir:}
  \lnot(!A \lif !B) & \Entails !A \land \lnot !B\\
  !A \land \lnot !B & \Entails \lnot(!A \lif !B)
  \intertext{Maddi koşul modus ponens'i destekler:}
  !A, !A \lif !B & \Entails !B
  \intertext{Maddi koşul birleştirilebilir:}
  !A \lif !B,  !A \lif !C & \Entails !A \lif (!B \land !C)
  \intertext{Önbileşeni her zaman güçlendirebiliriz; yani koşul
    \emph{monotondur}:}
  !A \lif !B & \Entails (!A \land !C) \lif !B
  \intertext{Maddi koşul geçişlidir; yani zincir kuralı geçerlidir:}
  !A \lif !B, !B \lif !C & \Entails !A \lif !C
  \intertext{Maddi koşul karşıt tersine denktir:}
  !A \lif !B & \Entails \lnot !B \lif \lnot !A\\
  \lnot !B \lif \lnot !A & \Entails !A \lif !B
\end{align}

Bunların tümü matematiksel akıl yürütmede yararlı ve sorunsuz çıkarımlardır.
Ne var ki felsefe ve dilbilim yazını, matematik dışındaki bağlamlarda bunlara
karşılık gelen çıkarımlara yöneltilen sözde karşı örneklerle doludur. Bu
örnekler, maddi koşul~$\lif$'in İngilizcedeki ``if \dots{} then \dots''
ifadelerini simgeleştirmek için uygun bağlaç olmadığını---ya da en azından
her zaman uygun olmadığını---düşündürür.

